\documentclass[11pt]{article}
\usepackage{fontspec}
\usepackage[french]{babel}
\usepackage{amsmath}
\usepackage{amssymb}
\usepackage{amsthm}
\usepackage{mathrsfs}
\usepackage{tikz-cd}
\usepackage{enumitem}

\newcommand{\printedpage}[1]{\marginpar{\footnotesize p.~#1}}
\newcommand{\msp}[1]{\marginpar{\footnotesize [#1]}}
\newcommand{\Hom}{\operatorname{Hom}}
\newcommand{\Ob}{\operatorname{Ob}}
\newcommand{\Fl}{\operatorname{Fl}}
\newcommand{\Hot}{\mathrm{Hot}}
\newcommand{\Cat}{\mathrm{Cat}}
\newcommand{\Ens}{\mathrm{Ens}}
\newcommand{\id}{\mathrm{id}}

% footnote marks in the text printed as superscript (n), in the footnote as superscript n
\makeatletter
\renewcommand\@makefnmark{\hbox{\@textsuperscript{\normalfont(\@thefnmark)}}}
\renewcommand\@makefntext[1]{\parindent 1em\noindent\hb@xt@1.8em{\hss\@textsuperscript{\normalfont\@thefnmark}}#1}
\makeatother

\begin{document}

\printedpage{9}
% head: A. Grothendieck : Dérivateurs, Ch. XVI, Édition des Universités de Montpellier II et Paris VII
% foot: Version 18 mars 2019

\[
\begin{tikzcd}[column sep=small]
\Hom(X,I_{\mathcal{H}}) \arrow[rr] \arrow[dr,dotted,"\pi_0^{X,I_{\mathcal{H}}}"'] & & \Hom(\pi_0(X),I) \arrow[dl,dotted,"\sim"'] \\ % ?? uncertain: label of the dotted right arrow, a tilde printed along the arrow -- \sim
& \Hom(\pi_0(X),\pi_0(I_{\mathcal{H}})). &
\end{tikzcd}
\]

La “bi-adjonction” des foncteurs $\pi_0$ et $I\longmapsto I_{\mathcal{H}}$ revient à dire que cette flèche est bijective,
\msp{17}
ou encore, que la flèche $\pi_0^{X,I_{\mathcal{H}}}$ :
\[
\Hom(X,I_{\mathcal{H}})\longrightarrow\Hom(\Gamma_{\mathcal{H}}(X),\Gamma_{\mathcal{H}}(I_{\mathcal{H}})\simeq I)
\]
soit bijective.

\textbf{Corollaire 2.11.} \textit{Le foncteur} $I\longmapsto I_W$ \textit{de} Ens \textit{dans} $\Hot_W$ \textit{est pleinement fidèle,} i.e. \textit{l'homomorphisme d'adjonction}
\[
\pi_0(I_W)\longrightarrow I
\]
\textit{est bijectif.} (\textit{Où encore,} % [sic] Où encore -> Ou encore, the accent makes no sense here
\textit{$\pi_0$ est un foncteur le localisation.}) % [sic] le localisation -> de localisation, wrong word

En effet, cette bijectivité est évidente quand on travaille dans Cat, et s'en déduit aussitôt dans $\Hot_W$.

\textsc{Démonstration de} 2.8. La formule d'adjonction (2.8.2) résulte de la formule correspondante pour Cat, évidente, et du lemme 2.6 sur les foncteurs adjoints entre catégories de modèles. L'adjonction dans l'autre sens se réduit à la formule (2.8.4). On a des flèches canoniques
\[
\begin{tikzcd}
\Hom_W(e_W,X) \arrow[r,shift left,"\alpha"] & \pi_0(X), \arrow[l,shift left,"\beta"]
\end{tikzcd}
\]
où $\alpha$ est $\pi_0^{e_W,X}$, compte tenu que $\pi_0(e_W)=\pi_0(e_{\Cat})=e_{\Ens}$, et $\Hom(e_{\Ens},\pi_0(X))\simeq\pi_0(X)$. La flèche $\beta$ est la flèche composée
\[
\Ob X\xrightarrow{\ \sim\ }\Hom_{\Cat}(e_{\Cat},X)\xrightarrow{\mathrm{hot}_W}\Hom_W(e_W,X)
\]
\msp{18}
en notant que cette flèche passe au quotient en $\pi_0(X)\longrightarrow\Hom_W(e_W,X)$, vu que deux applications homotopes de $e_{\Cat}$ dans $X$ sont égales dans $\Hot_W$. Je dis que
\[
\beta\alpha=\id\ ,\qquad\qquad \alpha\beta=\id\ .
\]
La relation $\alpha\beta=\id$ revient à dire que la flèche $i_x : e_{\Cat}\longrightarrow X$ définie par $x\in\Ob X$, donne $\pi_0(i_x) : e_{\Ens}\longrightarrow\pi_0(X)$ dont la valeur est la classe de $x$ dans $\pi_0(X)$, c'est tautologique. Pour prouver que $\beta\alpha=\id$, il suffit alors de prouver que $\beta$ est surjective, \textit{i.e.} le

\textbf{Corollaire 2.12.} \textit{Pour tout $X$ dans} Cat, \textit{l'application}
\[
\underbrace{\Hom_{\Cat}(e_{\Cat},X)}_{\simeq\Ob X}\longrightarrow\Hom_W(e_W,X)
\]
\printedpage{10}
% head: A. Grothendieck : Dérivateurs, Ch. XVI, Édition des Universités de Montpellier II et Paris VII
% foot: Version 18 mars 2019
\textit{est} surjective (\textit{moyennant l'axiome de connexité} Loc (5)).

\textsc{Démonstration de} 2.12. Si $X=\emptyset$, l'assertion revient à dire que $\Hom_W(e_W,\emptyset)=\emptyset$, mais plus généralement, on a le

\textbf{Corollaire 2.13.} (Moyennant Loc (5).) \textit{Soit $X$ non vide dans} Cat, \textit{alors} $\Hom_W(X,\emptyset)=\emptyset$. (\textit{Donc $\emptyset$ est un objet initial} strict \textit{de} $\Hot_W$.)

En effet, si on avait $X\longrightarrow\emptyset$ dans $\Hot_W$,
\msp{page 13'}\texttt{[il y a un décalage dans la numérotation]}
on en déduirait $\pi_0(X)\longrightarrow\emptyset$, ou $\pi_0(X)=\emptyset$, absurde\,!

Reste le cas où $X$ est non vide. Cela résulte du fait général :

\textbf{Lemme 2.14.} \textit{Soit $\mathcal{M}$ une catégorie ayant un objet initial} strict (\textit{si} $\Hom_{\mathcal{M}}(X,\emptyset_{\mathcal{M}})\neq\emptyset$, \textit{on a} $X=\emptyset_{\mathcal{M}}$), \textit{un objet final $e$, et telle que pour tout $X\neq\emptyset_{\mathcal{M}}$, $X\longrightarrow e$ ait une section,} i.e. $\Hom(e,X)\neq\emptyset$. \textit{Soit} $W\subset\Fl\mathcal{M}$
\begin{itemize}[leftmargin=3em,labelwidth=2.5em,align=right]
\item[1$^{\circ}$)] \textit{ne contenant pas de flèches} $\emptyset_{\mathcal{M}}\longrightarrow X$ \textit{avec} $X\neq\emptyset_{\mathcal{M}}$, \textit{et}
\item[2\,$^{\circ}$)] \textit{tel que pour tout diagramme}
\begin{tikzcd}[baseline=(\tikzcdmatrixname-1-1.base)]
X_1 & X_2 \arrow[l,"u_2"'] \\
e \arrow[u,"u_1"] \arrow[ur,dotted,"s"'] &
\end{tikzcd}
\textit{avec} $u_2\in W$ $\exists$ $s$ dans $\mathcal{M}$ \textit{tel que} $u_2s=u_1$

\textit{dans} $\mathcal{M}W^{-1}$.
\end{itemize}
\textit{Alors}
\begin{itemize}[leftmargin=3em,labelwidth=2.5em,align=right]
\item[a)] $\emptyset_{\mathcal{M}}$ \textit{est un objet initial} strict \textit{de} $\mathcal{M}W^{-1}$.
\item[b)] \textit{Pour tout} $X\neq\emptyset_{\mathcal{M}}$, $\Hom_{\mathcal{M}}(e,X)\longrightarrow\Hom_W(e,X)$ \textit{est surjectif.}
\end{itemize}

\textsc{Démonstration}. Que $\emptyset_{\mathcal{M}}$ soit objet initial de $\mathcal{M}W^{-1}$ a été vu (dernière assertion du lemme 2.5, forme duale), prouvons qu'il est strict, \textit{i.e.} que si $\Hom_W(X,\emptyset_{\mathcal{M}})\neq\emptyset$, on a $X=\emptyset_{\mathcal{M}}$. En effet, si on avait $u : X\longrightarrow\emptyset_{\mathcal{M}}$ dans $\mathcal{M}W^{-1}$, $u$ serait représenté par
\[
X=X_0\xleftarrow{\ W\ }X_1\longrightarrow X_2\xleftarrow{\ W\ }\cdots X_{2n-2}\longrightarrow X_{2n-1}\xleftarrow{\ W\ }X_{2n}=\emptyset_{\mathcal{M}}.
\]
On procède par récurrence sur la demi-longueur de cette chaîne pour voir que $X=\emptyset_{\mathcal{M}}$. En effet, l'hypothèse sur $W$ implique que $X_{2n-1}=\emptyset_{\mathcal{M}}$, puis celle que $\emptyset_{\mathcal{M}}$ est strict implique
$X_{2n-1}=\emptyset_{\mathcal{M}}$, d'où la conclusion. % [sic] X_{2n-1} -> X_{2n-2}, the strictness of the initial object gives X_{2n-2}

\msp{page 14'}
Prouvons la surjectivité de $\Hom_{\mathcal{M}}(e,X)\longrightarrow\Hom_W(e,X)$ pour $X\neq\emptyset_{\mathcal{M}}$, en montrant que tout chemin réduit de longueur 2n+1 est “homotope” à un chemin réduit de longueur $2n-1$ :
\[
e=X_0\longrightarrow X_1\xleftarrow{\ W\ }X_2\longrightarrow X_3\longleftarrow\ \cdots\ \longleftarrow X_{2n}\longrightarrow X_{2n+1}=X.
\]
En effet, on peut remplacer la section
\[
\begin{tikzcd}
e=X_0 \arrow[r,"u_1"'] \arrow[rr,dotted,bend right,"s"'] & X_1 & X_2 \arrow[l,"W"',"u_2"] \arrow[r,"u_3"'] & X_3 & X_4 \arrow[l,"W"']
\end{tikzcd}
\]
\printedpage{11}
% head: A. Grothendieck : Dérivateurs, Ch. XVI, Édition des Universités de Montpellier II et Paris VII
% foot: Version 18 mars 2019
par
\[
\begin{tikzcd}
e \arrow[r,"u_3s"'] & X_3 & \cdots, \arrow[l,"W"']
\end{tikzcd}
\]
où $s$ est une section de $X_2$ sur $e$ telle que
\[
u_2s=u_1\qquad\quad\textit{dans}\qquad\quad\mathcal{M}W^{-1}\ .
\]
Cela prouve donc par récurrence sur $n$ que le chemin est homotope à un chemin de longueur 1, $e\longrightarrow X$, q.e.d.

Il faut donc, dans le cas actuel, vérifier 1$^{\circ}$) (c'est essentiellement le corollaire 2.12.), et 2$^{\circ}$)
\[
\begin{tikzcd}
X_1 & X_2 \arrow[l,"u_2\in W"'] & \\
e \arrow[u,"u_1"] \arrow[ur,dotted,"s"'] & & .
\end{tikzcd}
\]
Mais comme $u_2\in W$, on a $\pi_0(u_2)$ isomorphisme, donc
\msp{page 15'}
surjectif, donc $\exists$ $s$ tel que $u_2s$ soit homotope à $u_1$ dans Cat, ce qui implique $u_2s=u_1$ dans $\Hot_W$, q.e.d.

\textbf{Commentaire final.} L'axiome de connexité Loc (5) diffère de tous les autres axiomes Loc (1) à Loc (8) (mais à l'exception de Loc (8 bis) d'accessibilité) par le fait que ce n'est pas un axiome de stabilité. Les autres axiomes Loc (1) à Loc (8) (et aussi Loc (8 bis) d'accessibilité et même les axiomes ultérieures % [sic] ultérieures -> ultérieurs, agreement with axiomes
d'exactitude~\footnote[7]{il n'y en aura en fait pas besoin, grâce à la providentielle théorie de K. S. \textsc{Brown}\,!}) sont satisfaits si $W=\Fl\mathcal{M}$ - l'axiome de connexité est donc le seul qui s'oppose à la \textit{trivialité} de la théorie homotopique définie par $W$. C'est dire qu'il ne peut être conséquence de l'un des autres. Par contre, l'axiome Loc (4) des sommes directes n'est pas indépendant des autres. Signalons à ce propos (rappel) :
\[
\left\{
\begin{array}{rl}
\text{Loc (4 a$_0$)} & \text{est conséquence de l'axiome Loc (6 a) du carré cocartésien (forme directe).}\\[1ex]
\text{Loc (4 a)} & \text{est conséquence de Loc (6 a) précédent}\\
 & \text{+ l'axiome des limites inductives filtrantes Loc (8).}\\[3ex]
\text{Loc (4 b)} & \text{est conséquence de l'axiome Loc (7) des fibrations. Donc}\\[1ex]
\text{Loc (4)} & \text{est conséquence de Loc (6 a, 7, 8).}
\end{array}
\right.
\]

\printedpage{12}
% head: A. Grothendieck : Dérivateurs, Ch. XVI, Édition des Universités de Montpellier II et Paris VII
% foot: Version 18 mars 2019
\msp{page 16'}

\section*{3\quad Axiome Loc (6) du carré cocartésien (ou de “Mayer-Vietoris”, ou des $W$-cofibrations)~\footnote[8]{C'est le nom “\textit{axiome des $W$-cofibrations}” qui m'apparaît à présent comme le meilleur, -- c'est ce nom qui me paraît le mieux saisir la signification conceptuelle et non technique de l'axiome.}}

Je vois plusieurs variantes, et n'ai pas encore l'esprit totalement net sur les implications mutuelles. Je commence par la variante plus compliquée, mais dont on aura besoin en tous cas :

\begin{itemize}[leftmargin=5em,labelwidth=4em,align=left]
\item[\textbf{Loc 6}] (Axiome des $W$-cofibrations, ou lemme des 5, ou axiome des homomorphismes de carrés cocartésiens). Considérons dans Cat un diagramme commutatif cubique
\[
\begin{tikzcd}[row sep=small,column sep=small]
X_0 \arrow[rr] \arrow[dr] \arrow[dd,"f_0"'] & & X_2 \arrow[dr] \arrow[dd,"f_2" near start] & \\
& X_1 \arrow[rr,crossing over] & & X_3 \arrow[dd,"f_3"] \\
Y_0 \arrow[rr] \arrow[dr] & & Y_2 \arrow[dr] & \\
& Y_1 \arrow[rr] \arrow[from=uu,crossing over,"f_1"' near start] & & Y_3,
\end{tikzcd}
\]
\end{itemize}
homomorphisme $f=(f_0,f_1,f_2,f_3)$ d'un diagramme carré $X$ dans un autre $Y$. On suppose que $X$ provient d'un objet $X_3$ de Cat de deux ouverts $X_1$, $X_2$ de $X_3$ \textit{recouvrant} $X_3$, et de leur intersection $X_0=X_1\cap X_2$, et de même pour $Y$. (Appelons ces diagrammes dans Cat des diagrammes MV (“\textit{ouverts}”) - “\textit{Mayer-Vietoris}”.) On a alors ce qui suit~\footnote[9]{Il est bon d'introduire des propriétés du diagramme
\begin{itemize}[leftmargin=3em,labelwidth=2.5em,align=right]
\item[1$^{\circ}$)] $f_0\in W$.
\item[2$^{\circ}$)] $f_1$ et $f_2$ dans $W$.
\item[3$^{\circ}$)] $f_3$ dans $W$.
\end{itemize}
et de formuler les parties a) b) c) de Loc (6) ainsi
\begin{itemize}[leftmargin=3em,labelwidth=2.5em,align=right]
\item[a)] 1$^{\circ}$ + 2$^{\circ}$ $\Longrightarrow$ 3$^{\circ}$.
\item[b)] 2$^{\circ}$ + 3$^{\circ}$ $\Longrightarrow$ 1$^{\circ}$.
\item[c)] 1$^{\circ}$ + 3$^{\circ}$ $\Longrightarrow$ 2$^{\circ}$.
\end{itemize}} :
\begin{itemize}[leftmargin=3em,labelwidth=2.5em,align=right]
\item[a)] $f_0,f_1,f_2\in W\Longrightarrow f_3\in W$~\footnote[10]{N.B. il semble que dans Loc (6 a), il suffit de l'énoncé en supposant $f_0$, $f_2$ des \textit{isomorphismes,} \textit{cf.} p. 22.}.
\item[b)] $f_1,f_2,f_3\in W\Longrightarrow f_0\in W$.
\item[c)] $f_0,f_3\in W\Longrightarrow f_1,f_2\in W$~\footnote[11]{Il ne faudrait pas inclure c dans Loc 6 \textit{cf.} p. 26.}.
\end{itemize}
\noindent\footnote[12]{N.B. On verra plus bas (3.3) que Loc 6 c implique déjà Loc (6 a, b), \textit{i.e.} Loc (6) tout entier. Et si $W$ satisfait Loc (1 bis), \textit{i.e.} stable par facteurs directs, alors Loc (6 a, b) $\Longrightarrow$ Loc (6 c).}.

\printedpage{13}
% head: A. Grothendieck : Dérivateurs, Ch. XVI, Édition des Universités de Montpellier II et Paris VII
% foot: Version 18 mars 2019
\msp{page 17'}
\noindent\footnote[13]{Je vais désigner par Loc (6) l'axiome Loc (6 a, b), sans y inclure c.}.

\textbf{N.B.} Je distingue ces conditions comme Loc (6 a), Loc (6 b), Loc (6 c). Je présume qu'on n'aura jamais à faire usage à Loc (6 b), sans supposer également Loc (6 a), mais l'inverse n'est sûrement pas le cas. Ainsi, $W_0=\{f\in\Fl\Cat\mid\pi_0(f)\text{ bijectif}\}$ satisfait a), mais non b). Je présume (sans avoir vérifié) que c'est pareil pour les $W_n$, $n\in\mathbf{N}$, de la “\textit{$n$-équivalence homotopique}”.

Il se pourrait que Loc (6 a) et Loc (6 a, b) soient équivalentes à la forme plus faible, et plus simple, suivante. (Si oui, c'est sous la forme la plus simple, bien sûr, qu'il convient de les énoncer.)

\begin{itemize}[leftmargin=5em,labelwidth=4em,align=left]
\item[\textbf{Loc 6$'$}] (Axiome des carrés cocartésiens, ou de Mayer-Vietoris). Soit
\[
\begin{tikzcd}[row sep=small,column sep=small]
X_0 \arrow[rr,"i_2"] \arrow[dr,"i_1"] & & X_2 \arrow[dr,"j_1"] & \\
& X_1 \arrow[rr,"j_2"] & & X_3
\end{tikzcd}
\]
un diagramme MV dans Cat. Alors
\begin{itemize}[leftmargin=3em,labelwidth=2.5em,align=right]
\item[a)] Si $i_1\in W$, $j_1\in W$.
\item[b)] Si $j_1\in W$, $i_1\in W$.
\end{itemize}
\end{itemize}

\msp{page 18'}
\textbf{Proposition 3.1.}
\[
\begin{aligned}
\text{Loc (6 a)}\quad &\Longrightarrow\quad \text{Loc (6$'$ a),}\\
\text{Loc (6 b)}\quad &\Longrightarrow\quad \text{Loc (6$'$ b),}\\
\text{Loc (6 c)}\quad &\Longrightarrow\quad \text{Loc (6$'$ a, b)}\ \footnotemark[14].
\end{aligned}
\]
\footnotetext[14]{N.B. on a même Loc (6 c) $\Longrightarrow$ $\underbrace{\text{Loc 6 (a, b)}}_{\text{Loc (6)}}$.}

Ça se voit sur le diagramme cubique
\[
\text{(Q)}\qquad\qquad
\begin{tikzcd}[row sep=small,column sep=small]
X_0 \arrow[rr,"i_2"] \arrow[dr,"i_1"] & & X_2 \arrow[dr,"j_1"] \arrow[from=dd,"f_2=\id"' near start] & \\
& X_1 \arrow[rr,"j_2",crossing over] & & X_3 \\
X_0 \arrow[uu,"f_0=\id"] \arrow[rr,"i_2"] \arrow[dr,"\id"] & & X_2 \arrow[dr,"\id"] & \\
& X_0 \arrow[rr,"i_2"] \arrow[uu,crossing over,"f_1=i_1" near end] & & X_2, \arrow[uu,"f_3=j_1"']
\end{tikzcd}
\]
où les $f_0$, $f_1$, $f_2$, $f_3$ de l'énoncé de Loc (6) deviennent $\id_{X_0}$, $i_1$, $\id_{X_2}$, $j_1$. Loc (6 a) donne donc $i_1\in W\Longrightarrow j_1\in W$, \textit{i.e.} Loc (6$'$ a), et Loc (6 c) $j_1\in W\Longrightarrow i_1\in W$, \textit{i.e.} Loc (6$'$ b). Pour Loc (6 b) $\Longrightarrow$ Loc (6$'$ b), on utilise (\textit{cf.} XII, p. 78) le cube

\end{document}
